%%
%% For submission and review, use:
%%   \documentclass[manuscript, screen, review]{acmart}
%% For final camera-ready, use:
%%   \documentclass[sigconf]{acmart}
%%
\documentclass[manuscript, screen, review]{acmart}

\AtBeginDocument{%
  \providecommand\BibTeX{{Bib\TeX}}}

%% Rights — fill in from ACM rights confirmation email after acceptance
\setcopyright{acmlicensed}
\copyrightyear{2026}
\acmYear{2026}
\acmDOI{XXXXXXX.XXXXXXX}
\acmConference[AIFM '26]{2026 2nd International Conference on Artificial Intelligence and Foundation Model}{June 26--28, 2026}{Urumqi, China}
\acmISBN{979-8-4007-2458-9}

%%
%% Submission ID — uncomment and fill in when assigned
%%\acmSubmissionID{XXX-XXXX}

\begin{document}

%%
%% Title. Use \title[short title]{full title} when the full title is long.
%%
\title[RegDivergence-101]{RegDivergence-101: An LLM Benchmark for
  Cross-Jurisdiction Regulatory Contradiction Detection in Life Sciences}

%%
%% Authors
%%
\author{Chuchu Wu}
\authornote{Corresponding author.}
\email{chuchuwu@outlook.com}
\orcid{0009-0006-5140-2247}
\affiliation{%
  \institution{Carnegie Mellon University}
  \city{Pittsburgh}
  \state{PA}
  \country{USA}
}

\author{Zhiyin Zhou}
\email{zzhou20@pratt.edu}
\affiliation{%
  \institution{Pratt Institute}
  \city{New York}
  \state{NY}
  \country{USA}
}

\author{Jingzhuo Hu}
\email{jingzhuoh00@gmail.com}
\affiliation{%
  \institution{University of Pennsylvania}
  \city{Philadelphia}
  \state{PA}
  \country{USA}
}

\renewcommand{\shortauthors}{Wu et al.}

%%
%% Abstract
%%
\begin{abstract}
Pharmaceutical sponsors developing a drug for both the United States and the
European Union must reconcile guidance issued independently by the FDA and the
EMA. Where the two agencies require \textit{substantively the same thing}, a
sponsor can file once; where they \textit{diverge}, a single trial design risks
rejection in one region; where one agency is \textit{silent} on a point the
other regulates, the sponsor must infer obligations. Today this reconciliation
is performed manually by regulatory-affairs experts. We introduce
\textbf{cross-jurisdiction regulatory divergence detection}: given an FDA
requirement and an EMA requirement on the same topic, classify their
relationship as AGREE, DIVERGE, or SILENT. We release
\textbf{RegDivergence-101}, a 101-pair expert-grounded benchmark (labels
sourced from three peer-reviewed FDA/EMA comparison studies; inter-annotator
$\kappa = 0.54$), and systematically
characterize a four-method baseline hierarchy: lexical heuristic (0.549
macro-F1, 95\% CI [0.458--0.650]), NLI cross-encoder (0.271),
obligation-level Graph-RAG (0.698 [0.602--0.783]), and flat LLM judge /
Claude Haiku (0.819 [0.736--0.900]). Three findings emerge: SILENT is
semantically detectable but invisible to entailment-only formulations;
pair-level obligation graphs improve over lexical methods but lose to flat-LLM
context; and corpus-level graph construction is the correct architectural
target for large-scale silent-detection.
\end{abstract}

%%
%% CCS Concepts — generated from https://dl.acm.org/ccs
%%
\begin{CCSXML}
<ccs2012>
  <concept>
    <concept_id>10010147.10010257</concept_id>
    <concept_desc>Computing methodologies~Natural language processing</concept_desc>
    <concept_significance>500</concept_significance>
  </concept>
  <concept>
    <concept_id>10002951.10003317.10003347</concept_id>
    <concept_desc>Information systems~Document filtering</concept_desc>
    <concept_significance>300</concept_significance>
  </concept>
  <concept>
    <concept_id>10003456.10003457.10003527</concept_id>
    <concept_desc>Applied computing~Health informatics</concept_desc>
    <concept_significance>300</concept_significance>
  </concept>
</ccs2012>
\end{CCSXML}

\ccsdesc[500]{Computing methodologies~Natural language processing}
\ccsdesc[300]{Information systems~Document filtering}
\ccsdesc[300]{Applied computing~Health informatics}

%%
%% Keywords
%%
\keywords{large language models, LLM benchmark, regulatory NLP, AI for
  compliance, FDA/EMA divergence, contradiction detection, Graph-RAG,
  retrieval-augmented generation, generative AI, life sciences,
  AGREE/DIVERGE/SILENT classification}

\maketitle

%% -------------------------------------------------------------------------
\section{Introduction}
%% -------------------------------------------------------------------------

Bringing a medicine to market in both the US and EU requires satisfying two
regulators that publish guidance separately and, frequently,
\textit{differently}. The canonical example is the choice of primary
\textbf{estimand} under ICH E9(R1): for the same trial, FDA guidance favors the
\textit{treatment-policy} strategy while EMA guidance favors the
\textit{hypothetical} strategy. Other documented divergences span pediatric
extrapolation (ICH E11A), endpoint definitions in ulcerative colitis, gene-therapy
long-term follow-up duration, and the regulatory pathway for companion
diagnostics. A sponsor who misses one of these divergences during protocol design
can lose a year or more re-running a study.

Reconciling the two corpora is currently a manual, expert task. Yet the NLP
building blocks for automating it all exist---document-level NLI
(ContractNLI~\cite{koreeda2021contractnli}, DocNLI~\cite{yin2021docnli}),
legal contradiction detection (LegalWiz~\cite{legalwiz2025},
CLAUSE~\cite{betterclause2026}), and Graph-RAG (Microsoft
GraphRAG~\cite{edge2024graphrag}, GraphCompliance~\cite{graphcompliance2025},
RAGulating Compliance~\cite{ragulating2025}). \textbf{No prior work combines
them to detect divergence between two parallel regulatory corpora from different
jurisdictions.} All existing contradiction work is \textit{intra-document}
(within one contract) or \textit{single-jurisdiction} (contract-vs-statute).

This paper makes three contributions:
\begin{enumerate}
\item \textbf{Task formulation.} Cross-jurisdiction regulatory divergence
  detection as a three-way classification (AGREE / DIVERGE / SILENT) over
  aligned FDA/EMA requirement pairs, grounded in published expert comparisons.
\item \textbf{Benchmark.} \textbf{RegDivergence-101}, 101 expert-sourced
  pairs across three therapeutic domains with inter-annotator agreement
  ($\kappa = 0.54$) and bootstrap confidence intervals.
\item \textbf{Baseline hierarchy with three findings.} Four methods
  systematically compared, yielding empirical insights on the task's structure
  and the path to a corpus-level solution.
\end{enumerate}

%% -------------------------------------------------------------------------
\section{Related Work}
%% -------------------------------------------------------------------------

\textbf{Document-level NLI.} ContractNLI~\cite{koreeda2021contractnli} introduced
document-level entailment/contradiction/not-mentioned classification with
evidence extraction, and showed that contradiction detection lags entailment,
especially under ``negation by exception.'' DocNLI~\cite{yin2021docnli} extended
NLI to full documents.

\textbf{Legal contradiction detection (2025--2026).}
LegalWiz~\cite{legalwiz2025} uses a multi-agent framework for legal
contradiction detection but is confined to single-document regimes. Better Call
CLAUSE~\cite{betterclause2026} (EACL 2026 Findings) benchmarks 7,500+
perturbed contracts across ten anomaly categories and finds that LLMs miss
subtle errors and struggle to \textit{justify} them legally. Contradiction
Detection in RAG Systems~\cite{contradictionrag2025} uses LLMs as context
validators. Two recent RAG-faithfulness studies bear directly on our LLM-judge
and graph-RAG results: \citet{chen2026ragwrong} probe whether a RAG system can
recognize when retrieved context is incorrect, and \citet{qian2026evidenceforce}
show that the surface relevance of retrieved evidence does not guarantee it
\textit{warrants} a claim---the same failure mode behind our Finding~2, where
modal-register similarity (RECOMMENDED vs.\ MANDATORY) misleads the pair-level
graph.

\textbf{Graph-RAG and regulatory compliance.}
Microsoft GraphRAG~\cite{edge2024graphrag} builds entity graphs with community
summaries for local and global queries. GraphCompliance~\cite{graphcompliance2025}
aligns \textit{policy graphs} (regulatory requirements) with \textit{context
graphs} (organizational facts) for GDPR; it explicitly lists ``handling
conflicting regulatory interpretations'' as open. RAGulating
Compliance~\cite{ragulating2025} builds an ontology-free regulatory knowledge
graph for traceable QA.

\textbf{Gap.} Contradiction detection is intra-document or
single-jurisdiction; Graph-RAG targets single-corpus QA or
single-jurisdiction compliance; FDA/EMA divergence is documented manually but
never automated. Our task sits precisely in this gap.

%% -------------------------------------------------------------------------
\section{Task and Dataset}
%% -------------------------------------------------------------------------

\subsection{Task Definition}

For a topic $t$, let $f_t$ be an FDA requirement statement and $e_t$ the
corresponding EMA statement. The label $y_t$ is:
\begin{itemize}
\item \textbf{AGREE} --- both jurisdictions require substantively the same thing.
\item \textbf{DIVERGE} --- the requirements directly conflict (one mandates
  what the other forbids, qualifies, or sets a different threshold for).
\item \textbf{SILENT} --- one jurisdiction does not address what the other
  regulates.
\end{itemize}

\subsection{Dataset Construction}

We seed pairs from documented FDA/EMA divergences in three open-access
expert-panel reviews: a \textit{J.\ Crohn's \& Colitis} 2025 expert comparison
of FDA and EMA ulcerative colitis guidance (40 pairs)~\cite{uc2025jcc}, a PMC
study comparing EMA and FDA psychiatric drug-trial guidelines (10
pairs)~\cite{pmc8157504}, and a PMC comparison of US and EU
radiopharmaceutical nonclinical requirements (9 pairs)~\cite{pmc6529498};
supplemented by 42 pairs derived from primary guidance text (ICH E9(R1),
E11A~\cite{iche11a2024}; FDA and EMA guidance). Each pair records \texttt{\{id,
topic, source, fda\_text, ema\_text, label\}}. The combined release contains
\textbf{101 pairs}: 38 AGREE, 40 DIVERGE, 23 SILENT, across 13 topic areas.
We report macro-F1 to weight all classes equally.

\textbf{Inter-annotator agreement.} To validate the labeling schema, a second
independent annotator (Claude Haiku with a deliberately different prompt and no
access to the original labels) re-labeled a stratified sample of 36 pairs (12
per class). Observed agreement: 0.694; Cohen's $\kappa =$ \textbf{0.542}
(moderate, Landis \& Koch 1977~\cite{landis1977kappa}). AGREE was the most
reliable class (83\% agreement); SILENT the hardest (58\%, mainly confused with
DIVERGE). We disclose that both annotation passes are LLM-assisted; the
expert-panel source papers provide the citable human-expert anchor. The moderate
$\kappa$ is below the conventional 0.60 threshold and reflects genuine task
difficulty---SILENT requires determining \textit{absence} of regulation, which
is inherently more ambiguous than detecting presence of conflict. We treat
$\kappa = 0.542$ as a characterization of task difficulty rather than a quality
failure, and recommend that a full human-expert IAA study be conducted for the
camera-ready benchmark release.

%% -------------------------------------------------------------------------
\section{Baselines}
%% -------------------------------------------------------------------------

We evaluate four methods under a common evaluation protocol:
\textbf{stratified nested 5-fold cross-validation} (inner folds tune any
hyperparameters, outer folds report, averaged over 5 random seeds) with
\textbf{bootstrap 95\% CI} (1,000 resamples). The LLM judge and Graph-RAG
classifier have no tunable hyperparameters and are evaluated directly; their
CIs are bootstrap-only.

\subsection{Lexical Heuristic}

No neural model, no API---fully offline and deterministic. Combines TF-IDF
cosine similarity with five conflict signals: negation asymmetry, divergence
cue pairs (e.g., \textit{single-arm} vs.\ \textit{randomized}), modal-strength
asymmetry (\textit{must/shall} vs.\ \textit{may/encouraged}),
numeric-threshold mismatch, and one-sided specificity (concrete number
vs.\ case-by-case deferral). A shared ICH citation is a strong AGREE signal.

\subsection{NLI Cross-Encoder}

\texttt{typeform/distilbert-base-uncased-mnli}, run in both directions; labels
derived from contradiction/neutral/entailment probabilities with CV-tuned
thresholds. We use DistilBERT-MNLI as a representative off-the-shelf NLI
system. The structural SILENT gap finding (\S\ref{sec:discussion}, Finding~1)
is model-agnostic: any NLI model lacking an explicit SILENT output class will
map absence-of-regulation to \textit{neutral}, collapsing SILENT recall.
Substituting DeBERTa-v3-large or a modern cross-encoder would improve
AGREE/DIVERGE F1 but would not resolve the SILENT structural gap without
task-specific supervision.

\subsection{Graph-RAG Pair-Level Classifier}

For each text, Claude Haiku extracts a policy triplet
$\langle$\textit{subject, obligation\_level, requirement, conditions}$\rangle$
where $\text{obligation\_level} \in \{\text{MANDATORY}, \text{PROHIBITED},
\text{RECOMMENDED}, \text{PERMITTED}, \text{SILENT}\}$. Only
MANDATORY~$\leftrightarrow$~PROHIBITED is a hard-coded DIVERGE rule; all other
pairs are escalated to a second LLM call that reasons over the
\textit{obligation structure} of the aligned node pair, with explicit
calibration that modal-register differences (RECOMMENDED vs.\ MANDATORY) do not
automatically imply DIVERGE.

\textbf{Evaluation note.} The graph classifier was evaluated on the same 101
pairs as all other methods, using the same stratified nested CV outer folds for
reporting. No hyperparameter tuning was performed (the MANDATORY
$\leftrightarrow$ PROHIBITED rule and the LLM prompt were fixed before
evaluation began). Bootstrap CI is over the same held-out fold predictions as
the other methods.

\subsection{LLM Judge (Flat Pairwise)}

Claude Haiku with explicit AGREE / DIVERGE / SILENT definitions, prompted once
per pair with both texts. No graph, no retrieval. Full prompt in
Appendix~\ref{app:prompt}.

%% -------------------------------------------------------------------------
\section{Results}
%% -------------------------------------------------------------------------

Table~\ref{tab:results} reports macro-F1, accuracy, and per-class F1 with
bootstrap 95\% confidence intervals for all four methods.

\begin{table*}
  \caption{Macro-F1, accuracy, and per-class F1 with bootstrap 95\% confidence
    intervals for the four baselines ($n = 101$ pairs).}
  \label{tab:results}
  \begin{tabular}{llllll}
    \toprule
    Method & Macro-F1 [95\% CI] & Acc & AGREE F1 [CI] & DIVERGE F1 [CI] & SILENT F1 [CI] \\
    \midrule
    Lexical heuristic$^1$
      & 0.556 [0.458--0.650] & 0.586
      & 0.756 [0.648--0.849] & 0.538 [0.400--0.660] & 0.361 [0.167--0.545] \\
    NLI cross-encoder$^2$
      & 0.271 [0.181--0.354] & 0.277
      & 0.303 [0.118--0.481] & 0.416 [0.278--0.545] & 0.079 [0.000--0.177] \\
    Graph-RAG pair-level$^3$
      & 0.698 [0.602--0.783] & 0.703
      & 0.766 [0.646--0.860] & 0.691 [0.554--0.806] & 0.628 [0.468--0.764] \\
    \textbf{LLM judge (Haiku)}$^4$
      & \textbf{0.819 [0.736--0.900]} & \textbf{0.832}
      & \textbf{0.884 [0.800--0.951]} & \textbf{0.830 [0.737--0.911]}
      & \textbf{0.743 [0.579--0.878]} \\
    \bottomrule
  \end{tabular}
  \smallskip\\
  {\small $^1$Nested 5-fold CV, 5 seeds; thresholds tuned on inner folds only.
  $^2$Same CV protocol, thresholds tuned on inner folds.
  $^3$No tunable hyperparameters; bootstrap CI only.
  $^4$\texttt{claude-haiku-4-5-20251001}, default temperature; bootstrap CI
  only. Estimated cost: {\raise.17ex\hbox{$\scriptstyle\sim$}}\$0.10 for all
  101 pairs.}
\end{table*}

CIs do not overlap between lexical and Graph-RAG, or between Graph-RAG and LLM
judge---the gaps are statistically reliable at $n = 101$.

%% -------------------------------------------------------------------------
\section{Error Analysis: The Lexical Ceiling}
%% -------------------------------------------------------------------------

The seven residual lexical errors share a signature: \textbf{high lexical
overlap, opposed regulatory stance.}

\begin{itemize}
\item \textbf{adaptive-01 (DIVERGE $\to$ missed):} FDA ``does \textit{not
  require} detailed pre-specification'' vs.\ EMA ``should be \textit{fully
  prespecified}''---near-identical vocabulary, opposite obligation.
\item \textbf{uc-01:} Single-arm Mayo-score endpoint vs.\ co-primary
  modified-Mayo---same terms, different structural requirement.
\item \textbf{gene-01:} FDA recommends 15 years; EMA defers to
  case-by-case---same topic, different specificity.
\end{itemize}

A bag-of-words model cannot resolve these; they require representing
\textit{what each requirement asserts about which entity}.

%% -------------------------------------------------------------------------
\section{Discussion: What the Baseline Hierarchy Reveals}
\label{sec:discussion}
%% -------------------------------------------------------------------------

\textbf{Finding 1 --- SILENT is semantically detectable but invisible to NLI
framing.} SILENT F1 = 0.079 under NLI (no native SILENT class) vs.\ 0.743
under the LLM judge (explicit definition). This is not a model-capacity
effect---it is structural. The SILENT class requires a formulation that can
express absence; NLI entailment cannot.

\textbf{Finding 2 --- Pair-level obligation graphs improve over lexical (+0.14
F1) but lose to flat LLM ($-$0.12 F1).} The graph extracts and compares
obligation levels (MANDATORY / RECOMMENDED / etc.), which helps distinguish
clear stance differences from similar-sounding text. But it loses the
full-text context that the flat LLM uses to separate
\textit{same-requirement / different-modal-register} (AGREE) from
\textit{same-topic / different-threshold} (DIVERGE). Pairs like gene-01 (FDA
15 years vs.\ EMA case-by-case) both extract as RECOMMENDED and the graph
calls them AGREE; the flat LLM correctly identifies the threshold conflict from
the text. This is precisely the failure mode identified by
\citet{qian2026evidenceforce}: surface relevance does not imply evidential
force.

\textbf{Finding 3 --- Two distinct ceilings for two different problem scopes.}
The \textit{pair-level ceiling} is the flat LLM judge (0.819): given two
aligned sentences and explicit definitions, an LLM nearly saturates the task.
The \textit{corpus-level ceiling} is the domain of the graph method: ``is EMA
silent on this topic \textit{anywhere} across its guideline corpus?'' is a
query a pairwise LLM cannot make but a corpus-level policy graph can. These
are different tasks with different upper bounds. This paper measures and
establishes the pair-level ceiling; \S\ref{sec:futurework} motivates the
corpus-level architecture as the natural next contribution.

%% -------------------------------------------------------------------------
\section{Future Work: Corpus-Level Policy Graph}
\label{sec:futurework}
%% -------------------------------------------------------------------------

The pair-level experiments establish that the correct target for a genuine
Graph-RAG contribution is \textbf{corpus-level} silent detection and divergence
discovery---not re-ranking an already-aligned pair, but \textit{finding} which
FDA requirements have no EMA counterpart and vice versa.

The proposed architecture:
\begin{enumerate}
\item \textbf{Per-jurisdiction policy-graph construction}---index all FDA
  guidance documents and all EMA guidelines; extract requirement triplets at
  corpus scale.
\item \textbf{Cross-jurisdiction alignment}---entity normalization +
  embedding similarity to link corresponding requirement nodes across the two
  graphs.
\item \textbf{Divergence and silence discovery}---for each aligned node pair:
  DIVERGE if obligation structures conflict; SILENT if a concept node exists in
  one graph but has no aligned counterpart in the other.
\item \textbf{Explanation generation}---natural-language summaries for
  regulatory-affairs specialists.
\end{enumerate}

Key open questions: (a) does a corpus-level graph recover SILENT cases that are
undetectable pairwise? (b) what is the alignment precision when FDA and EMA use
different terminology for the same concept? (c) can the graph be kept current
as guidance is updated?

%% -------------------------------------------------------------------------
\section{Limitations}
%% -------------------------------------------------------------------------

\textbf{Benchmark scale.} 101 pairs from three source documents introduces
domain skew (40/101 from one UC comparison paper) and limits statistical power.
The CI widths ($\approx$0.09--0.17 for most methods) reflect this. Results
should be treated as directional.

\textbf{Inter-annotator agreement.} $\kappa = 0.542$ is moderate. Both
annotation passes are LLM-assisted; no independent human-expert double
annotation was performed. The expert-panel source papers provide the primary
validity anchor, but a full human IAA study is needed for a final benchmark
release.

\textbf{Data contamination.} Claude Haiku may have been exposed to FDA/EMA
guidance and the published expert-panel comparison papers during pretraining. To
probe this, we split performance by source: on the 42 \textit{primary guidance
pairs} (low contamination risk) Haiku achieves AGREE F1 = 0.941 and DIVERGE F1
= 0.968; on the 59 \textit{expert-comparison pairs} (higher risk) AGREE F1 =
0.769 and DIVERGE F1 = 0.745. Performance is \textit{not} higher on the
potentially-contaminated split. The macro-F1 gap (0.636 vs.\ 0.765) is fully
explained by SILENT support ($n = 2$ on primary vs.\ $n = 21$ on expert
pairs), not memorization. Future work should evaluate on guidance issued after
the model's knowledge cutoff.

\textbf{SILENT operationalization.} Our SILENT pairs use a placeholder on the
absent side, which may not reflect real retrieval conditions where a system must
determine absence from a full corpus. The pair-level SILENT task is a proxy;
\S\ref{sec:futurework} describes the correct corpus-level formulation.

%% -------------------------------------------------------------------------
\section{Conclusion}
%% -------------------------------------------------------------------------

We introduced cross-jurisdiction regulatory divergence detection, released
\textbf{RegDivergence-101}, a 101-pair expert-grounded benchmark
($\kappa = 0.54$), and characterized a
four-method baseline hierarchy with bootstrap CIs. Three empirical findings:
SILENT requires explicit absence-aware formulation (NLI F1 0.08 $\to$ LLM F1
0.74); obligation-level graph structure improves over lexical heuristics (+0.14
F1) but loses to flat LLM context ($-$0.12 F1); and the graph method's correct
target is corpus-level construction, where it uniquely enables silence
discovery at scale. The flat LLM judge (0.819) is the pair-level ceiling;
beating it requires the corpus-level architecture proposed in
\S\ref{sec:futurework}.

%%
%% Acknowledgments
%%
\begin{acks}
This work was conducted independently. Benchmark pairs were derived from
publicly available FDA/EMA guidance documents and open-access peer-reviewed
comparison studies. LLM API calls used the Anthropic API (Claude Haiku,
\texttt{claude-haiku-4-5-20251001}).
\end{acks}

%%
%% Bibliography
%%
\bibliographystyle{ACM-Reference-Format}
\bibliography{xjurisdiction-divergence}

%%
%% Appendix
%%
\appendix

\section{LLM Judge Prompt}
\label{app:prompt}

\textbf{System:}
\begin{verbatim}
You are a regulatory-affairs expert comparing FDA and EMA
guidance documents. Given an FDA requirement and an EMA
requirement on the same topic, classify their relationship.

Respond with ONLY a JSON object in this exact format:
{"label": "<AGREE|DIVERGE|SILENT>", "reason": "<one sentence>"}

Definitions:
- AGREE: both jurisdictions require substantively the same thing
- DIVERGE: the requirements directly conflict (one mandates
  what the other forbids or softens)
- SILENT: one jurisdiction does not address what the other
  regulates
\end{verbatim}

\textbf{User:}
\begin{verbatim}
FDA requirement:
{fda_text}

EMA requirement:
{ema_text}

Classify the relationship.
\end{verbatim}

Model: \texttt{claude-haiku-4-5-20251001}. Temperature: default (1.0).
Max tokens: 100. Results cached; each pair scored exactly once.

\end{document}
